\chapter{Regular Curves}
\label{ch:vii12-regular-curves}

A parametrized curve is the simplest immersed manifold. Its geometry begins with a distinction that is easy to overlook: the image of a curve and its parametrization are not the same object. Regularity makes the velocity nonzero, while reparametrization isolates those quantities that belong to the geometric curve rather than to the chosen clock.

\section{Parametrized and regular curves}
\begin{definition}[Parametrized curve]\label{def:vii12-param-curve}
A smooth parametrized curve in a smooth manifold $M$ is a smooth map
\[
\gamma:I\to M
\]
from an interval $I\subset\mathbb R$.
\end{definition}

\begin{definition}[Regular curve]\label{def:vii12-regular}
A parametrized curve is regular when
\[
\dot\gamma(t)=d\gamma_t\left(\frac d{dt}\right)\ne0
\]
for every $t\in I$. Equivalently, $\gamma$ is an immersion.
\end{definition}

\begin{example}[Regular and singular planar curves]\label{ex:vii12-regular-singular}
The circle $(\cos t,\sin t)$ is regular. The cusp $(t^2,t^3)$ is not regular at $t=0$ because its velocity vanishes there.
\end{example}

\begin{proposition}[Local embeddedness]\label{prop:vii12-local-embedded}
Every regular curve is locally an embedded one-dimensional submanifold of $M$.
\end{proposition}
\begin{proof}
A regular curve has rank one. The constant-rank theorem gives local coordinates in which it is $t\mapsto(t,0,\ldots,0)$ after a local change of parameter.
\end{proof}

\section{Reparametrization}
\begin{definition}[Regular reparametrization]\label{def:vii12-reparam}
A reparametrization is a diffeomorphism $\phi:J\to I$. The same geometric curve is then represented by
\[
\widetilde\gamma=\gamma\circ\phi.
\]
It is orientation preserving when $\phi'>0$ and reversing when $\phi'<0$.
\end{definition}

\begin{proposition}[Velocity transformation]\label{prop:vii12-velocity-transform}
Under reparametrization,
\[
\dot{\widetilde\gamma}(s)=\dot\gamma(\phi(s))\,\phi'(s).
\]
Hence regularity is invariant under reparametrization.
\end{proposition}

\section{Curves in Euclidean space}
\begin{definition}[Speed]\label{def:vii12-speed}
For a smooth curve $\gamma:I\to\mathbb R^n$, its speed is
\[
v(t)=\|\dot\gamma(t)\|.
\]
A regular curve has strictly positive speed.
\end{definition}

\begin{definition}[Length]\label{def:vii12-length}
For a smooth curve on $[a,b]$ define
\[
L(\gamma)=\int_a^b\|\dot\gamma(t)\|\,dt.
\]
For piecewise smooth curves, sum the lengths of the smooth pieces.
\end{definition}

\begin{theorem}[Length is reparametrization invariant]\label{thm:vii12-length-invariant}
If $\phi:[c,d]\to[a,b]$ is a diffeomorphism, then
\[
L(\gamma\circ\phi)=L(\gamma).
\]
\end{theorem}
\begin{proof}
Use $\|\dot\gamma(\phi(s))\phi'(s)\|=\|\dot\gamma(\phi(s))\|\,|\phi'(s)|$ and the one-variable change-of-variables theorem.
\end{proof}

\section{Arc length}
\begin{definition}[Arc-length function]\label{def:vii12-arclength-function}
For a regular curve $\gamma:[a,b]\to\mathbb R^n$, define
\[
s(t)=\int_a^t\|\dot\gamma(u)\|\,du.
\]
Then $s'(t)=\|\dot\gamma(t)\|>0$.
\end{definition}

\begin{theorem}[Unit-speed reparametrization]\label{thm:vii12-unit-speed}
Every regular Euclidean curve admits a local, and on an interval global, orientation-preserving reparametrization by arc length. In the new parameter $s$,
\[
\left\|\frac{d\gamma}{ds}\right\|=1.
\]
\end{theorem}
\begin{proof}
Since $s'(t)>0$, the inverse-function theorem gives an inverse $t=t(s)$. Then
\[
\left\|\frac{d\gamma}{ds}\right\|=\|\dot\gamma(t(s))\|\frac{dt}{ds}=1.
\]
\end{proof}

\begin{definition}[Unit tangent]\label{def:vii12-unit-tangent}
For a regular Euclidean curve,
\[
T(t)=\frac{\dot\gamma(t)}{\|\dot\gamma(t)\|}
\]
is the unit tangent. Under orientation-preserving reparametrization it is unchanged; under orientation reversal it changes sign.
\end{definition}

\section{Simple and closed curves}
\begin{definition}[Simple curve]\label{def:vii12-simple}
A parametrized curve is simple when it is injective, except that for a closed curve one permits the two endpoints of a compact parameter interval to have the same image.
\end{definition}

\begin{definition}[Closed curve]\label{def:vii12-closed}
A smooth curve $\gamma:[a,b]\to M$ is smoothly closed when $\gamma(a)=\gamma(b)$ and the endpoint derivatives match under the identification needed to descend to a smooth map $S^1\to M$.
\end{definition}

\begin{warning}[Same image, different traversal]\label{warning:vii12-image-param}
Two curves may have the same image but traverse it with different speed, orientation, or multiplicity. Geometric invariants must therefore be checked for reparametrization invariance.
\end{warning}

\section{Standard examples}
\begin{example}[Line]\label{ex:vii12-line}
For $\gamma(t)=p+tv$ with $v\ne0$, speed is $\|v\|$ and length on $[a,b]$ is $(b-a)\|v\|$.
\end{example}

\begin{example}[Circle]\label{ex:vii12-circle}
For $\gamma(t)=(R\cos t,R\sin t)$, speed is $R$ and one full traversal has length $2\pi R$.
\end{example}

\begin{example}[Helix]\label{ex:vii12-helix}
For $\gamma(t)=(R\cos t,R\sin t,ct)$,
\[
\|\dot\gamma(t)\|=\sqrt{R^2+c^2},
\]
so length over $[a,b]$ is $(b-a)\sqrt{R^2+c^2}$.
\end{example}

\section{Curves on manifolds}
\begin{proposition}[Coordinate velocity]\label{prop:vii12-coordinate-velocity}
If $x=(x^1,\ldots,x^n)$ is a chart and $\gamma(t)$ lies in its domain, then
\[
\dot\gamma(t)=\sum_i\frac{d}{dt}(x^i\circ\gamma)(t)\left.\frac{\partial}{\partial x^i}\right|_{\gamma(t)}.
\]
\end{proposition}

\begin{proposition}[Velocity under smooth maps]\label{prop:vii12-map-velocity}
For a smooth map $F:M\to N$,
\[
\frac d{dt}(F\circ\gamma)(t)=dF_{\gamma(t)}\dot\gamma(t).
\]
\end{proposition}

\section{Piecewise regular curves and line integrals}
\begin{definition}[Piecewise regular]\label{def:vii12-piecewise}
A continuous curve on $[a,b]$ is piecewise regular if there is a finite subdivision on whose subintervals it is smooth and regular, with one-sided derivatives at break points.
\end{definition}

\begin{definition}[Line integral of a one-form]\label{def:vii12-line-integral}
For a piecewise smooth curve $\gamma:[a,b]\to M$ and a one-form $\alpha$ defined along its image,
\[
\int_\gamma\alpha:=\int_a^b\alpha_{\gamma(t)}(\dot\gamma(t))\,dt.
\]
This is invariant under orientation-preserving reparametrization and changes sign when orientation is reversed.
\end{definition}

\begin{proposition}[Exact one-forms along curves]\label{prop:vii12-exact-line}
For a smooth function $f$,
\[
\int_\gamma df=f(\gamma(b))-f(\gamma(a)).
\]
\end{proposition}
\begin{proof}
$df_{\gamma(t)}(\dot\gamma(t))=\frac d{dt}(f\circ\gamma)(t)$, so integrate and apply the fundamental theorem of calculus.
\end{proof}

\section*{Bridge to VII/13}
Regularity and arc length provide the correct parameter-independent starting point for local curve geometry. VII/13 differentiates the unit tangent, constructs Frenet frames, and introduces curvature and torsion; those quantities are intentionally deferred until the parametrization issues have been settled here.


\section{Exercises}

\begin{exercise}\label{exr:vii12-01}
Determine whether $\gamma(t)=(t,t^2)$ is regular.
\end{exercise}
\begin{hint}
% PEDAGOGY-ENRICHED-VII
Differentiate. Make the controlling geometric criterion explicit before the calculation. Parametrize by an allowed regular parameter, compute speed and arc length, and separate geometric quantities from artifacts of reparametrization.
\end{hint}
\begin{solution}
$\dot\gamma=(1,2t)$ never vanishes, so the curve is regular.
\end{solution}

\begin{exercise}\label{exr:vii12-02}
Determine where $\gamma(t)=(t^2,t^3)$ fails to be regular.
\end{exercise}
\begin{hint}
% PEDAGOGY-ENRICHED-VII
Differentiate. Work in a convenient local model first, then return to the invariant statement. Parametrize by an allowed regular parameter, compute speed and arc length, and separate geometric quantities from artifacts of reparametrization.
\end{hint}
\begin{solution}
$\dot\gamma=(2t,3t^2)$ vanishes exactly at $t=0$.
\end{solution}

\begin{exercise}\label{exr:vii12-03}
Show a regular curve is an immersion.
\end{exercise}
\begin{hint}
% PEDAGOGY-ENRICHED-VII
The domain is one-dimensional. After the main computation, check the hypotheses and geometric interpretation. Parametrize by an allowed regular parameter, compute speed and arc length, and separate geometric quantities from artifacts of reparametrization.
\end{hint}
\begin{solution}
Its differential has rank one exactly when the velocity is nonzero, which is the definition of regularity.
\end{solution}

\begin{exercise}\label{exr:vii12-04}
Explain why a regular curve is locally embedded.
\end{exercise}
\begin{hint}
% PEDAGOGY-ENRICHED-VII
Use the constant-rank theorem. Make the controlling geometric criterion explicit before the calculation. Parametrize by an allowed regular parameter, compute speed and arc length, and separate geometric quantities from artifacts of reparametrization.
\end{hint}
\begin{solution}
Rank one gives local coordinates in which the map is $t\mapsto(t,0,\ldots,0)$, an embedding onto a coordinate line.
\end{solution}

\begin{exercise}\label{exr:vii12-05}
For $\widetilde\gamma=\gamma\circ\phi$, compute its velocity.
\end{exercise}
\begin{hint}
% PEDAGOGY-ENRICHED-VII
Apply the chain rule. Work in a convenient local model first, then return to the invariant statement. Parametrize by an allowed regular parameter, compute speed and arc length, and separate geometric quantities from artifacts of reparametrization.
\end{hint}
\begin{solution}
$\dot{\widetilde\gamma}(s)=\dot\gamma(\phi(s))\phi'(s)$.
\end{solution}

\begin{exercise}\label{exr:vii12-06}
Why does a diffeomorphic reparametrization preserve regularity?
\end{exercise}
\begin{hint}
% PEDAGOGY-ENRICHED-VII
Its derivative never vanishes. After the main computation, check the hypotheses and geometric interpretation. Parametrize by an allowed regular parameter, compute speed and arc length, and separate geometric quantities from artifacts of reparametrization.
\end{hint}
\begin{solution}
Both factors in $\dot\gamma(\phi(s))\phi'(s)$ are nonzero for a regular curve and a parameter diffeomorphism.
\end{solution}

\begin{exercise}\label{exr:vii12-07}
Compute the speed of $\gamma(t)=(3t,4t)$.
\end{exercise}
\begin{hint}
% PEDAGOGY-ENRICHED-VII
Take the Euclidean norm. Make the controlling geometric criterion explicit before the calculation. Parametrize by an allowed regular parameter, compute speed and arc length, and separate geometric quantities from artifacts of reparametrization.
\end{hint}
\begin{solution}
The speed is $\sqrt{3^2+4^2}=5$.
\end{solution}

\begin{exercise}\label{exr:vii12-08}
Compute the length of $\gamma(t)=(3t,4t)$ on $[0,2]$.
\end{exercise}
\begin{hint}
% PEDAGOGY-ENRICHED-VII
Integrate its constant speed. Work in a convenient local model first, then return to the invariant statement. Parametrize by an allowed regular parameter, compute speed and arc length, and separate geometric quantities from artifacts of reparametrization.
\end{hint}
\begin{solution}
$L=\int_0^2 5\,dt=10$.
\end{solution}

\begin{exercise}\label{exr:vii12-09}
Prove length invariance under orientation-preserving reparametrization.
\end{exercise}
\begin{hint}
% PEDAGOGY-ENRICHED-VII
Substitute $t=\phi(s)$. After the main computation, check the hypotheses and geometric interpretation. Parametrize by an allowed regular parameter, compute speed and arc length, and separate geometric quantities from artifacts of reparametrization.
\end{hint}
\begin{solution}
The transformed speed is $\|\dot\gamma(\phi(s))\|\phi'(s)$, and ordinary substitution yields the original length integral.
\end{solution}

\begin{exercise}\label{exr:vii12-10}
Why is the absolute value of $\phi'$ needed for orientation-reversing reparametrizations?
\end{exercise}
\begin{hint}
% PEDAGOGY-ENRICHED-VII
Length is nonnegative. Make the controlling geometric criterion explicit before the calculation. Parametrize by an allowed regular parameter, compute speed and arc length, and separate geometric quantities from artifacts of reparametrization.
\end{hint}
\begin{solution}
The norm gives $\|\dot\gamma(\phi(s))\phi'(s)\|=\|\dot\gamma(\phi(s))\||\phi'(s)|$, so reversal does not make length negative.
\end{solution}

\begin{exercise}\label{exr:vii12-11}
Compute the speed of the radius-$R$ circle $(R\cos t,R\sin t)$.
\end{exercise}
\begin{hint}
% PEDAGOGY-ENRICHED-VII
Differentiate. Work in a convenient local model first, then return to the invariant statement. Parametrize by an allowed regular parameter, compute speed and arc length, and separate geometric quantities from artifacts of reparametrization.
\end{hint}
\begin{solution}
The velocity is $(-R\sin t,R\cos t)$, whose norm is $R$.
\end{solution}

\begin{exercise}\label{exr:vii12-12}
Compute one full circle length from its parametrization.
\end{exercise}
\begin{hint}
% PEDAGOGY-ENRICHED-VII
Integrate from $0$ to $2\pi$. After the main computation, check the hypotheses and geometric interpretation. Parametrize by an allowed regular parameter, compute speed and arc length, and separate geometric quantities from artifacts of reparametrization.
\end{hint}
\begin{solution}
$L=\int_0^{2\pi}R\,dt=2\pi R$.
\end{solution}

\begin{exercise}\label{exr:vii12-13}
Compute the speed of the helix $(R\cos t,R\sin t,ct)$.
\end{exercise}
\begin{hint}
% PEDAGOGY-ENRICHED-VII
Differentiate and square components. Make the controlling geometric criterion explicit before the calculation. Parametrize by an allowed regular parameter, compute speed and arc length, and separate geometric quantities from artifacts of reparametrization.
\end{hint}
\begin{solution}
The squared speed is $R^2\sin^2t+R^2\cos^2t+c^2=R^2+c^2$, so speed is $\sqrt{R^2+c^2}$.
\end{solution}

\begin{exercise}\label{exr:vii12-14}
Show the arc-length function of a regular curve is strictly increasing.
\end{exercise}
\begin{hint}
% PEDAGOGY-ENRICHED-VII
Differentiate it. Work in a convenient local model first, then return to the invariant statement. Parametrize by an allowed regular parameter, compute speed and arc length, and separate geometric quantities from artifacts of reparametrization.
\end{hint}
\begin{solution}
$s'(t)=\|\dot\gamma(t)\|>0$, so $s$ is strictly increasing.
\end{solution}

\begin{exercise}\label{exr:vii12-15}
Why can arc length be used as a parameter?
\end{exercise}
\begin{hint}
% PEDAGOGY-ENRICHED-VII
Apply the inverse-function theorem. After the main computation, check the hypotheses and geometric interpretation. Parametrize by an allowed regular parameter, compute speed and arc length, and separate geometric quantities from artifacts of reparametrization.
\end{hint}
\begin{solution}
Since $s'(t)>0$, $s$ is locally a diffeomorphism; on an interval it is globally one-to-one onto its image.
\end{solution}

\begin{exercise}\label{exr:vii12-16}
Show that an arc-length parametrized curve has unit speed.
\end{exercise}
\begin{hint}
% PEDAGOGY-ENRICHED-VII
Differentiate through the inverse $t(s)$. Make the controlling geometric criterion explicit before the calculation. Parametrize by an allowed regular parameter, compute speed and arc length, and separate geometric quantities from artifacts of reparametrization.
\end{hint}
\begin{solution}
$d\gamma/ds=\dot\gamma(t(s))/s'(t(s))$, whose norm is $1$.
\end{solution}

\begin{exercise}\label{exr:vii12-17}
How does the unit tangent transform under an orientation-preserving reparametrization?
\end{exercise}
\begin{hint}
% PEDAGOGY-ENRICHED-VII
Normalize the transformed velocity. Work in a convenient local model first, then return to the invariant statement. Parametrize by an allowed regular parameter, compute speed and arc length, and separate geometric quantities from artifacts of reparametrization.
\end{hint}
\begin{solution}
The positive factor $\phi'$ cancels in normalization, so $\widetilde T(s)=T(\phi(s))$.
\end{solution}

\begin{exercise}\label{exr:vii12-18}
How does the unit tangent transform under orientation reversal?
\end{exercise}
\begin{hint}
% PEDAGOGY-ENRICHED-VII
Now $\phi'<0$. After the main computation, check the hypotheses and geometric interpretation. Parametrize by an allowed regular parameter, compute speed and arc length, and separate geometric quantities from artifacts of reparametrization.
\end{hint}
\begin{solution}
Normalization leaves the sign of $\phi'$, so $\widetilde T(s)=-T(\phi(s))$.
\end{solution}

\begin{exercise}\label{exr:vii12-19}
Give an example of two parametrizations with the same image but different multiplicity.
\end{exercise}
\begin{hint}
% PEDAGOGY-ENRICHED-VII
Use a circle. Make the controlling geometric criterion explicit before the calculation. Parametrize by an allowed regular parameter, compute speed and arc length, and separate geometric quantities from artifacts of reparametrization.
\end{hint}
\begin{solution}
$(\cos t,\sin t)$ on $[0,2\pi]$ traverses once, while $(\cos2t,\sin2t)$ on the same interval traverses twice.
\end{solution}

\begin{exercise}\label{exr:vii12-20}
State the coordinate formula for velocity on a manifold.
\end{exercise}
\begin{hint}
% PEDAGOGY-ENRICHED-VII
Differentiate coordinate functions. Work in a convenient local model first, then return to the invariant statement. Parametrize by an allowed regular parameter, compute speed and arc length, and separate geometric quantities from artifacts of reparametrization.
\end{hint}
\begin{solution}
$\dot\gamma=\sum_i(d/dt)(x^i\circ\gamma)\,\partial/\partial x^i$.
\end{solution}

\begin{exercise}\label{exr:vii12-21}
Prove velocity is natural under smooth maps.
\end{exercise}
\begin{hint}
% PEDAGOGY-ENRICHED-VII
Use the chain rule for differentials. After the main computation, check the hypotheses and geometric interpretation. Parametrize by an allowed regular parameter, compute speed and arc length, and separate geometric quantities from artifacts of reparametrization.
\end{hint}
\begin{solution}
$d(F\circ\gamma)_t=dF_{\gamma(t)}\circ d\gamma_t$; applying to $d/dt$ gives the formula.
\end{solution}

\begin{exercise}\label{exr:vii12-22}
Define a piecewise regular curve.
\end{exercise}
\begin{hint}
% PEDAGOGY-ENRICHED-VII
Use a finite subdivision. Make the controlling geometric criterion explicit before the calculation. Parametrize by an allowed regular parameter, compute speed and arc length, and separate geometric quantities from artifacts of reparametrization.
\end{hint}
\begin{solution}
It is continuous and admits finitely many subintervals on which it is smooth and regular, with one-sided derivatives at the subdivision points.
\end{solution}

\begin{exercise}\label{exr:vii12-23}
Show the line integral of a one-form is invariant under orientation-preserving reparametrization.
\end{exercise}
\begin{hint}
% PEDAGOGY-ENRICHED-VII
Transform both velocity and $dt$. Work in a convenient local model first, then return to the invariant statement. Parametrize by an allowed regular parameter, compute speed and arc length, and separate geometric quantities from artifacts of reparametrization.
\end{hint}
\begin{solution}
The factor $\phi'$ from the velocity combines with $ds$ exactly as in one-variable substitution, giving the same integral.
\end{solution}

\begin{exercise}\label{exr:vii12-24}
Prove $\int_\gamma df=f(\gamma(b))-f(\gamma(a))$.
\end{exercise}
\begin{hint}
% PEDAGOGY-ENRICHED-VII
Use the chain rule. After the main computation, check the hypotheses and geometric interpretation. Parametrize by an allowed regular parameter, compute speed and arc length, and separate geometric quantities from artifacts of reparametrization.
\end{hint}
\begin{solution}
The integrand is $df(\dot\gamma)=d(f\circ\gamma)/dt$; integrate and apply the fundamental theorem.
\end{solution}


\section{Solved problem dossiers}

\begin{problem}[Legacy regular-curve problem]\label{prob:vii12-01}
Show that nonvanishing velocity is equivalent to the immersion condition for a curve.
\end{problem}
\begin{solution}
Since $T_tI$ is one-dimensional, $d\gamma_t$ is injective exactly when $d\gamma_t(d/dt)=\dot\gamma(t)$ is nonzero.
\end{solution}

\begin{problem}[Legacy parametrization problem]\label{prob:vii12-02}
Prove the reparametrization velocity formula and deduce invariance of regularity.
\end{problem}
\begin{solution}
Apply the chain rule to $\gamma\circ\phi$. Because a parameter diffeomorphism has $\phi'\ne0$, nonzero velocity is preserved.
\end{solution}

\begin{problem}[Legacy length problem]\label{prob:vii12-03}
Prove Euclidean curve length is invariant under all diffeomorphic reparametrizations.
\end{problem}
\begin{solution}
The transformed speed is the old speed times $|\phi'|$. The one-variable change-of-variables theorem, with orientation handled by the absolute value, recovers the same integral.
\end{solution}

\begin{problem}[Legacy arc-length problem]\label{prob:vii12-04}
Construct the arc-length parameter for a regular curve and prove it gives unit speed.
\end{problem}
\begin{solution}
Set $s(t)=\int_a^t\|\dot\gamma(u)\|du$. Since $s'>0$, invert to $t(s)$ and differentiate $\gamma(t(s))$ to obtain norm one.
\end{solution}

\begin{problem}[Line and segment]\label{prob:vii12-05}
Compute speed and length of $\gamma(t)=p+tv$ and explain its dependence on the parameter interval.
\end{problem}
\begin{solution}
Velocity is the constant vector $v$, so speed is $\|v\|$ and length on $[a,b]$ is $(b-a)\|v\|$. Reparametrizing changes the clock but not the geometric segment length.
\end{solution}

\begin{problem}[Circle]\label{prob:vii12-06}
Compute the length and unit tangent of a radius-$R$ circle.
\end{problem}
\begin{solution}
For $\gamma(t)=(R\cos t,R\sin t)$, speed is $R$, length over one turn is $2\pi R$, and $T(t)=(-\sin t,\cos t)$.
\end{solution}

\begin{problem}[Helix]\label{prob:vii12-07}
Compute speed and length of a circular helix.
\end{problem}
\begin{solution}
$\dot\gamma=(-R\sin t,R\cos t,c)$ has constant norm $\sqrt{R^2+c^2}$, so length is parameter-span times that constant.
\end{solution}

\begin{problem}[Local embeddedness]\label{prob:vii12-08}
Use the constant-rank theorem to show every regular curve is locally a coordinate axis.
\end{problem}
\begin{solution}
Rank one allows source and target coordinates in which the map is $s\mapsto(s,0,\ldots,0)$. Therefore the image near each point is an embedded one-dimensional submanifold.
\end{solution}

\begin{problem}[Same image, opposite orientation]\label{prob:vii12-09}
Compare $\gamma(t)=(\cos t,\sin t)$ with $\widetilde\gamma(t)=(\cos t,-\sin t)$.
\end{problem}
\begin{solution}
They have the same circle image and equal speed, but traverse it in opposite orientations. Their unit tangents and line integrals of one-forms differ by sign after the corresponding orientation reversal.
\end{solution}

\begin{problem}[Coordinate velocity on a manifold]\label{prob:vii12-10}
Derive the coordinate expression for $\dot\gamma$ from the derivation definition of tangent vectors.
\end{problem}
\begin{solution}
Apply $\dot\gamma$ to a smooth function and use the coordinate chain rule. The coefficients multiplying the coordinate derivations are the derivatives of the coordinate functions along the curve.
\end{solution}

\begin{problem}[Line integral of an exact form]\label{prob:vii12-11}
Show line integrals of $df$ depend only on endpoints.
\end{problem}
\begin{solution}
Pull $df$ back to the interval: $\gamma^*df=d(f\circ\gamma)$. Integration gives $f(\gamma(b))-f(\gamma(a))$.
\end{solution}

\begin{problem}[Piecewise regular concatenation]\label{prob:vii12-12}
Show line integrals and lengths add under concatenation of piecewise regular curves with matching endpoints.
\end{problem}
\begin{solution}
Choose a subdivision at the joining time. Both definitions are sums of ordinary integrals over the pieces, so additivity is immediate.
\end{solution}


\section{Challenge problems}

\begin{problem}[Arc length on a Riemannian manifold preview]\label{prob:vii12-ch01}
Assume a Riemannian metric $g$ is given. Show that $L_g(\gamma)=\int\sqrt{g(\dot\gamma,\dot\gamma)}dt$ is reparametrization invariant.
\end{problem}
\begin{solution}
Under reparametrization, the velocity is multiplied by $\phi'$. Bilinearity makes its squared norm multiply by $(\phi')^2$, so the speed multiplies by $|\phi'|$. One-variable substitution then proves invariance.
\end{solution}

\begin{problem}[Regular closed curve as an immersion of $S^1$]\label{prob:vii12-ch02}
Show that a smoothly closed regular curve with matching endpoint jets descends to a smooth immersion $S^1\to M$.
\end{problem}
\begin{solution}
Identify $S^1$ with $[a,b]/(a\sim b)$. Matching derivatives ensure smoothness across the quotient point, and nonzero velocity gives rank one everywhere.
\end{solution}

\begin{problem}[Length versus multiplicity]\label{prob:vii12-ch03}
If a regular closed curve is traversed $k$ times by an orientation-preserving $k$-fold parameter map, compare its length with one traversal.
\end{problem}
\begin{solution}
A $k$-fold traversal is not a diffeomorphic reparametrization of the once-traversed interval; it changes multiplicity. Splitting into $k$ monotone pieces shows the total length is $k$ times the one-traversal length.
\end{solution}

\begin{problem}[Line integrals detect orientation]\label{prob:vii12-ch04}
Let $\alpha=-y\,dx+x\,dy$ on the unit circle. Compare its integral under counterclockwise and clockwise parametrizations.
\end{problem}
\begin{solution}
Counterclockwise parametrization gives $\alpha(\dot\gamma)=1$ and integral $2\pi$. Reversing orientation multiplies the line integral by $-1$, giving $-2\pi$.
\end{solution}

\begin{problem}[Why curvature belongs in the next chapter]\label{prob:vii12-ch05}
Explain why $\|dT/dt\|$ is not itself a geometric invariant under arbitrary reparametrization, and identify the invariant combination that motivates arc length.
\end{problem}
\begin{solution}
Under an orientation-preserving change $t=t(s)$, $dT/ds=(dT/dt)(dt/ds)$, so its norm scales with parameter speed. Differentiating with respect to arc length removes this factor. Thus $\|dT/ds\|$ is the natural next invariant, which is precisely the curvature introduced in VII/13.
\end{solution}
